% Lecture 7: determinant properties and applications, subspaces, dimension, and rank.
% Sources: Lecture7/04.ppt (Sections 3.2--3.3) and Lecture7/05.pptx.
% Build from the repository root after copying this file to Lecture7:
%   pdflatex -interaction=nonstopmode -halt-on-error Lecture7/lecture7.tex
\pdfminorversion=7
\documentclass[aspectratio=169]{beamer}

\usepackage{graphicx}
\usepackage{amsmath,amssymb,mathtools}
\usepackage{xspace}
\usepackage{tikz}
\usepackage{xcolor}
\usepackage{colortbl}
\usepackage{booktabs}
\usepackage{array}
\usepackage{multirow}
\usepackage{pifont}
\usepackage{etoolbox}
\usepackage{arydshln}
\usepackage{iftex}
\ifPDFTeX
  \usepackage[UTF8,heading=false]{ctex}
\else
  \usepackage{fontspec}
  \usepackage{xeCJK}
  \setCJKmainfont{Microsoft YaHei}
\fi
\usetikzlibrary{arrows.meta,calc,fit,matrix,positioning,shapes.geometric}

\makeatletter
\def\input@path{{../}{Lecture7/}}
\makeatother

\setbeamercolor{headline}{fg=black,bg=black}
\setbeamercolor{section in head/foot}{bg=gray}
\setbeamertemplate{headline}{%
  \begin{beamercolorbox}[ht=2.25ex,dp=3.75ex]{section in head/foot}
    \insertnavigation{\paperwidth}
  \end{beamercolorbox}}

\mode<presentation>{
  \usetheme{OSUCSE169}
  \setbeamercovered{transparent=28}
}
\setbeamertemplate{footline}[frame number]
\setbeamertemplate{navigation symbols}{}
\setbeamercovered{invisible}

\definecolor{lecturegreen}{RGB}{48,105,21}
\definecolor{lectureblue}{RGB}{0,76,170}
\definecolor{lectureorange}{RGB}{202,105,0}
\definecolor{lecturemagenta}{RGB}{210,0,135}
\definecolor{theorembg}{RGB}{245,245,255}
\newcommand{\R}{\mathbb{R}}
\newcommand{\vect}[1]{\mathbf{#1}}
\newcommand{\Col}{\operatorname{Col}}
\newcommand{\Nul}{\operatorname{Nul}}
\newcommand{\rank}{\operatorname{rank}}
\newcommand{\Span}{\operatorname{Span}}
\newcommand{\sectioncard}[2]{%
  \begin{frame}[plain]
    \vfill\centering
    {\Huge\color{lectureorange}\bfseries #1\\[1ex]#2}
    \vfill
  \end{frame}}
\newcommand{\theorembox}[2]{%
  \begin{beamercolorbox}[sep=1.2ex,wd=\linewidth]{block body}
    \textbf{#1}\par #2
  \end{beamercolorbox}}

\begin{document}

\begin{frame}[plain]
  \vfill
  {\Huge\color{lecturegreen}\bfseries Lecture 7}
  \bigskip

  {\LARGE\color{lecturegreen}\bfseries Determinant Applications, Subspaces, Dimension, and Rank}
  \vfill
\end{frame}

\begin{frame}{Roadmap}
  \begin{enumerate}
    \item 3.2 Properties of Determinants
    \item 3.3 Applications of Determinant Properties
    \item 2.8 Subspaces of $\R^n$
    \item 2.9 Dimension and Rank
  \end{enumerate}
\end{frame}

\begin{frame}{Learning Goals}
  By the end of the lecture, you should be able to
  \begin{itemize}
    \item track determinant changes under row and column operations;
    \item connect determinants with invertibility and linear systems;
    \item apply Cramer's rule to small square systems;
    \item recognize subspaces and find bases for column and null spaces;
    \item compute coordinates, dimension, rank, and nullity.
  \end{itemize}
\end{frame}

% -----------------------------------------------------------------------------
\section[Determinant Properties]{Properties of Determinants}
\sectioncard{3.2}{Properties of Determinants}

\begin{frame}{Elementary Row Operations and Determinants}
Let $B$ be obtained from a square matrix $A$ by one elementary row operation.
\renewcommand{\arraystretch}{1.35}
\begin{center}
\begin{tabular}{>{\raggedright\arraybackslash}p{.47\linewidth}c}
\toprule
Operation from $A$ to $B$ & Effect on the determinant \\
\midrule
$R_i\leftrightarrow R_j$ & $\det B=-\det A$ \\
$R_i\leftarrow kR_i$ & $\det B=k\det A$ \\
$R_i\leftarrow R_i+cR_j$, $i\ne j$ & $\det B=\det A$ \\
\bottomrule
\end{tabular}
\end{center}
\end{frame}

\begin{frame}{A Row Interchange Reverses the Sign}
If two rows of $A$ are interchanged to produce $B$, then
\[
\boxed{\det B=-\det A}.
\]
For example,
\[
\begin{vmatrix}1&7&5\\6&6&2\\3&5&8\end{vmatrix}
=-
\begin{vmatrix}1&7&5\\3&5&8\\6&6&2\end{vmatrix}.
\]
\pause
After $r$ row interchanges, the accumulated factor is $(-1)^r$.
\end{frame}

\begin{frame}{Equal Rows Force a Zero Determinant}
Suppose $A$ has two equal rows. Interchanging them leaves $A$ unchanged, but the interchange rule gives
\[
\det A=-\det A.
\]
Therefore
\[
\boxed{\det A=0}.
\]
\pause
The same argument applies when two columns are equal.
\end{frame}

\begin{frame}{Scaling One Row}
If one row of $A$ is multiplied by $k$ to produce $B$, then
\[
\boxed{\det B=k\det A}.
\]
Only one row is scaled in this statement.
\pause

If every row of an $n\times n$ matrix is scaled by $k$, then
\[
\boxed{\det(kA)=k^n\det A}.
\]
In particular,
\[
\det(-A)=(-1)^n\det A.
\]
\end{frame}

\begin{frame}{Two Immediate Zero Tests}
\begin{block}{A zero row or column}
If $A$ contains a zero row or zero column, then $\det A=0$.
\end{block}
\pause
\begin{block}{A multiple row or column}
If one row is a scalar multiple of another row, then $\det A=0$. The same statement holds for columns.
\end{block}
\pause
For example,
\[
\begin{vmatrix}2&-4&5\\1&2&4\\-2&4&-5\end{vmatrix}=0
\]
because the third row is $-1$ times the first row.
\end{frame}

\begin{frame}{Row Replacement Preserves the Determinant}
If a multiple of one row is added to another row, then
\[
\boxed{\det B=\det A}.
\]
Thus row replacement can create zeros without adding a correction factor.
\pause

The same rule holds for column replacement.
\end{frame}

\begin{frame}{Cofactor Expansion After Row Replacement}
Compute
\[
D=\begin{vmatrix}
1&3&-1&2\\
1&-5&3&-4\\
0&2&1&-1\\
-5&1&3&-3
\end{vmatrix}.
\]
Use $R_2\leftarrow R_2-R_1$ and $R_4\leftarrow R_4+5R_1$:
\[
D=\begin{vmatrix}
1&3&-1&2\\0&-8&4&-6\\0&2&1&-1\\0&16&-2&7
\end{vmatrix}.
\]
\pause
Expanding down the first column and reducing the remaining determinant gives
\[
\boxed{D=-40}.
\]
\end{frame}

\begin{frame}{A Determinant Ledger}
Reduce $A$ to an upper triangular matrix $U$ and record each operation.
\begin{itemize}
  \item Row replacement contributes no factor.
  \item Each row interchange contributes a factor $-1$.
  \item Scaling a row by $k$ multiplies the determinant by $k$.
\end{itemize}
\pause
If elimination uses only row replacements and $r$ interchanges, then
\[
\boxed{\det A=(-1)^r u_{11}u_{22}\cdots u_{nn}}.
\]
\end{frame}

\begin{frame}{Elimination Example}
Compute
\[
\det\begin{bmatrix}2&4&6\\5&6&7\\7&6&10\end{bmatrix}.
\]
Factor $2$ from the first row, then use row replacement:
\[
2\begin{vmatrix}1&2&3\\5&6&7\\7&6&10\end{vmatrix}
=2\begin{vmatrix}1&2&3\\0&-4&-8\\0&-8&-11\end{vmatrix}
=2(-4)\begin{vmatrix}1&2&3\\0&1&2\\0&-8&-11\end{vmatrix}.
\]
One more replacement produces a final pivot $5$, so
\[
\boxed{\det A=2(-4)(1)(5)=-40}.
\]
\end{frame}

\begin{frame}{Use Structure Before Expanding}
Let
\[
D=\begin{vmatrix}
3&2&2&2\\2&3&2&2\\2&2&3&2\\2&2&2&3
\end{vmatrix}.
\]
Add all rows to the first row, giving $(9,9,9,9)$, and factor out $9$.
Then subtract the first row from the other rows:
\[
D=9\begin{vmatrix}
1&1&1&1\\0&1&0&0\\0&0&1&0\\0&0&0&1
\end{vmatrix}=\boxed{9}.
\]
\end{frame}

\begin{frame}{Combining Expansion and Reduction}
Compute
\[
\begin{vmatrix}
1&2&3&4\\0&5&0&0\\2&7&6&10\\2&9&7&11
\end{vmatrix}.
\]
Expand down the second row:
\[
=5\begin{vmatrix}1&3&4\\2&6&10\\2&7&11\end{vmatrix}.
\]
Use $R_2\leftarrow R_2-2R_1$ and $R_3\leftarrow R_3-2R_1$:
\[
5\begin{vmatrix}1&3&4\\0&0&2\\0&1&3\end{vmatrix}
=-5(1)(1)(2)=\boxed{-10}.
\]
\end{frame}

\begin{frame}{Transpose Property}
\theorembox{Theorem}{For every square matrix $A$,}
\[
\boxed{\det(A^T)=\det A}.
\]
\pause
Rows and columns therefore play symmetric roles. Every rule for row operations has a corresponding rule for column operations.
\end{frame}

\begin{frame}{Column Operations}
The determinant responds to column operations exactly as it responds to row operations:
\begin{itemize}
  \item interchanging two columns reverses the sign;
  \item multiplying one column by $k$ multiplies the determinant by $k$;
  \item adding a multiple of one column to another preserves the determinant.
\end{itemize}
\pause
A column operation on $A$ is the corresponding row operation on $A^T$.
\end{frame}

\begin{frame}{Multiplicative Property}
\theorembox{Theorem}{If $A$ and $B$ are $n\times n$ matrices, then}
\[
\boxed{\det(AB)=(\det A)(\det B)}.
\]
\pause
Consequently,
\[
\det(A^m)=(\det A)^m.
\]
For example, if $\det A=5$, then
\[
\det(A^3)=5^3=125.
\]
\end{frame}

\begin{frame}{Properties Summary}
For square matrices $A$ and $B$ and a scalar $c$:
\[
\det(AB)=\det A\det B,
\qquad
\det(A^T)=\det A,
\qquad
\det(cA)=c^n\det A.
\]
Also, $\det A=0$ if $A$ has
\begin{itemize}
  \item a zero row or column;
  \item two equal rows or columns;
  \item one row or column that is a multiple of another.
\end{itemize}
\end{frame}

% -----------------------------------------------------------------------------
\section[Determinant Applications]{Applications of Determinant Properties}
\sectioncard{3.3}{Applications of Determinant Properties}

\begin{frame}{Application 1: Efficient Computation}
Two ideas work well together:
\begin{enumerate}
  \item create zeros with row or column replacement;
  \item use cofactor expansion or finish with a triangular determinant.
\end{enumerate}
\pause
Before computing, check for
\begin{itemize}
  \item a row or column with many zeros;
  \item a quick row replacement that creates zeros;
  \item repeated or proportional rows or columns.
\end{itemize}
\end{frame}

\begin{frame}{Efficient Computation Example}
\small
Compute
\[
D=\begin{vmatrix}
1&4&5&0\\4&5&0&1\\5&0&1&4\\0&1&4&5
\end{vmatrix}.
\]
Adding the four rows to the first row gives $(10,10,10,10)$, so
\[
D=10\begin{vmatrix}
1&1&1&1\\4&5&0&1\\5&0&1&4\\0&1&4&5
\end{vmatrix}.
\]
Row replacement and a cofactor expansion reduce this to
\[
10\cdot 8\cdot(-8)\begin{vmatrix}3&2\\1&1\end{vmatrix}
=\boxed{-640}.
\]
\end{frame}

\begin{frame}{Application 2: The Invertibility Test}
\theorembox{Invertibility Theorem}{A square matrix $A$ is invertible if and only if $\det A\ne0$.}
\medskip
If $A$ is invertible, then
\[
1=\det I=\det(AA^{-1})=(\det A)(\det A^{-1}),
\]
so $\det A\ne0$.
\pause
Conversely, elimination shows that $\det A\ne0$ forces a pivot in every row and column, so $A$ is invertible.
\end{frame}

\begin{frame}{The Determinant of an Inverse}
If $A$ is invertible, then
\[
(\det A)(\det A^{-1})=\det I=1.
\]
Therefore
\[
\boxed{\det(A^{-1})=\frac{1}{\det A}}.
\]
\pause
If $AB=I$ for square matrices $A$ and $B$, then $\det A\det B=1$, so both matrices are invertible and
\[
B=A^{-1}.
\]
\end{frame}

\begin{frame}{A Singular Product}
Suppose $A$ and $B$ are $n\times n$ matrices, $\det B\ne0$, and $\det(AB)=0$.
Then
\[
0=\det(AB)=(\det A)(\det B).
\]
Because $\det B\ne0$, it follows that
\[
\boxed{\det A=0},
\]
so $A$ is singular.
\end{frame}

\begin{frame}{Equivalent Conditions for a Square Matrix}
For an $n\times n$ matrix $A$, the following statements are equivalent:
\begin{enumerate}
  \item $A$ is invertible.
  \item $A$ is row equivalent to $I_n$.
  \item $\det A\ne0$.
  \item $A\vect x=\vect b$ has a unique solution for every $\vect b\in\R^n$.
  \item $A\vect x=\vect0$ has only the trivial solution.
  \item $A^T$ is invertible.
  \item $\rank A=n$.
\end{enumerate}
\end{frame}

\begin{frame}{Application 3: Classifying Square Systems}
For a square system $A\vect x=\vect b$:
\begin{itemize}
  \item if $\det A\ne0$, the system has exactly one solution;
  \item if $\det A=0$, compare $\rank A$ and $\rank[A\mid\vect b]$.
\end{itemize}
\pause
When $\det A=0$,
\[
\begin{array}{ccl}
\rank A=\rank[A\mid\vect b]<n &\Longrightarrow& \text{infinitely many solutions},\\[1ex]
\rank A<\rank[A\mid\vect b] &\Longrightarrow& \text{no solution}.
\end{array}
\]
\end{frame}

\begin{frame}{A Parameterized System}
Classify the solutions of
\[
\left\{\begin{aligned}
kx+y+z&=1,\\
x+ky+z&=k,\\
x+y+kz&=k^2.
\end{aligned}\right.
\]
The coefficient matrix is
\[
A=\begin{bmatrix}k&1&1\\1&k&1\\1&1&k\end{bmatrix},
\qquad
\det A=(k+2)(k-1)^2.
\]
\end{frame}

\begin{frame}{The Parameterized System: Classification}
From $\det A=(k+2)(k-1)^2$:
\begin{itemize}
  \item if $k\ne1,-2$, then $\det A\ne0$ and the system has a unique solution;
  \item if $k=1$, then $\rank A=\rank[A\mid\vect b]=1<3$, so the system has infinitely many solutions;
  \item if $k=-2$, then $\rank A=2<\rank[A\mid\vect b]=3$, so the system has no solution.
\end{itemize}
\end{frame}

\begin{frame}{Homogeneous and Nonhomogeneous Systems}
For $A_{m\times n}\vect x=\vect0$:
\[
\begin{array}{ccl}
\rank A=n &\Longleftrightarrow& \text{only the trivial solution},\\
\rank A<n &\Longleftrightarrow& \text{nontrivial solutions exist}.
\end{array}
\]
For $A_{m\times n}\vect x=\vect b$:
\[
\begin{array}{ccl}
\rank A=\rank[A\mid\vect b]=n &\Longrightarrow& \text{one solution},\\
\rank A=\rank[A\mid\vect b]<n &\Longrightarrow& \text{infinitely many},\\
\rank A\ne\rank[A\mid\vect b] &\Longrightarrow& \text{no solution}.
\end{array}
\]
\end{frame}

\begin{frame}{Application 4: Cramer's Rule}
\theorembox{Cramer's Rule}{Let $A$ be an invertible $n\times n$ matrix. The unique solution of $A\vect x=\vect b$ is}
\[
\boxed{x_i=\frac{\det A_i(\vect b)}{\det A}},
\qquad i=1,2,\ldots,n,
\]
where $A_i(\vect b)$ is obtained from $A$ by replacing column $i$ with $\vect b$.
\pause
Cramer's rule requires $\det A\ne0$.
\end{frame}

\begin{frame}{Cramer's Rule Example}
Solve
\[
\left\{\begin{aligned}
x_1+2x_2-x_3&=-5,\\
2x_1-x_2+x_3&=6,\\
x_1-x_2-3x_3&=-3.
\end{aligned}\right.
\]
The determinants are
\[
\det A=19,
\qquad
\det A_1=19,
\qquad
\det A_2=-38,
\qquad
\det A_3=38.
\]
Thus
\[
\boxed{(x_1,x_2,x_3)=(1,-2,2)}.
\]
\end{frame}

\begin{frame}{A Four-Variable Example}
For the system
\[
\left\{\begin{aligned}
2x_1+x_2-5x_3+x_4&=8,\\
x_1-3x_2-6x_4&=9,\\
2x_2-x_3+2x_4&=-5,\\
x_1+4x_2-7x_3+6x_4&=0,
\end{aligned}\right.
\]
the source computation gives
\[
\det A=27,
\quad
(\det A_1,\det A_2,\det A_3,\det A_4)=(81,-108,-27,27).
\]
Therefore
\[
\boxed{(x_1,x_2,x_3,x_4)=(3,-4,-1,1)}.
\]
\end{frame}

\begin{frame}{When to Use Cramer's Rule}
Cramer's rule gives an exact formula and helps prove theoretical results.
\pause

For hand calculations, it is practical mainly for $2\times2$ and small $3\times3$ systems. Larger systems require many determinants, so elimination is usually more efficient.
\medskip

For a general linear system, use row reduction to determine consistency and solve. Use inverse matrices or Cramer's rule only when the coefficient matrix is square and invertible.
\end{frame}

% -----------------------------------------------------------------------------
\section[Subspaces]{Subspaces of Rn}
\sectioncard{2.8}{Subspaces of $\R^n$}

\begin{frame}{Span}
For vectors $\vect v_1,\ldots,\vect v_p\in\R^n$,
\[
\Span\{\vect v_1,\ldots,\vect v_p\}
=\left\{c_1\vect v_1+\cdots+c_p\vect v_p:c_1,\ldots,c_p\in\R\right\}.
\]
This set contains every linear combination of the given vectors.
\pause

We call it the subset of $\R^n$ spanned or generated by $\vect v_1,\ldots,\vect v_p$.
\end{frame}

\begin{frame}{Subspaces of \texorpdfstring{$\R^n$}{R^n}}
A subset $H$ of $\R^n$ is a \textbf{subspace} if
\begin{enumerate}
  \item $\vect0\in H$;
  \item $\vect u+\vect v\in H$ whenever $\vect u,\vect v\in H$;
  \item $c\vect u\in H$ whenever $\vect u\in H$ and $c\in\R$.
\end{enumerate}
\pause
The last two conditions say that $H$ is closed under addition and scalar multiplication.
\end{frame}

\begin{frame}{Every Span Is a Subspace}
Let $H=\Span\{\vect v_1,\vect v_2\}$. Then
\[
\vect0=0\vect v_1+0\vect v_2\in H.
\]
If
\[
\vect u=s_1\vect v_1+s_2\vect v_2,
\qquad
\vect v=t_1\vect v_1+t_2\vect v_2,
\]
then
\[
\vect u+\vect v=(s_1+t_1)\vect v_1+(s_2+t_2)\vect v_2\in H.
\]
Also, $c\vect u=(cs_1)\vect v_1+(cs_2)\vect v_2\in H$.
\end{frame}

\begin{frame}{Geometric Examples}
\begin{itemize}
  \item A line through the origin in $\R^2$ or $\R^3$ is a subspace.
  \item A plane through the origin in $\R^3$ is a subspace.
  \item $\R^n$ is a subspace of itself.
  \item The zero subspace $\{\vect0\}$ is a subspace.
\end{itemize}
\pause
\begin{alertblock}{The origin matters}
A line or plane that does not pass through the origin cannot be a subspace because it does not contain $\vect0$.
\end{alertblock}
\end{frame}

\begin{frame}{Column Space}
For $A=[\vect a_1\ \cdots\ \vect a_n]\in\R^{m\times n}$, the \textbf{column space} is
\[
\Col A=\Span\{\vect a_1,\ldots,\vect a_n\}.
\]
Thus $\Col A$ is a subspace of $\R^m$.
\pause

A vector $\vect b$ belongs to $\Col A$ exactly when
\[
A\vect x=\vect b
\]
is consistent.
\end{frame}

\begin{frame}{Column Space Example}
Let
\[
A=\begin{bmatrix}1&-3&-4\\-4&6&-2\\-3&7&6\end{bmatrix},
\qquad
\vect b=\begin{bmatrix}3\\3\\-4\end{bmatrix}.
\]
Row reduction gives
\[
\left[\begin{array}{ccc|c}
1&-3&-4&3\\-4&6&-2&3\\-3&7&6&-4
\end{array}\right]
\sim
\left[\begin{array}{ccc|c}
1&-3&-4&3\\0&-6&-18&15\\0&0&0&0
\end{array}\right].
\]
The system is consistent, so $\boxed{\vect b\in\Col A}$.
\end{frame}

\begin{frame}{Null Space}
The \textbf{null space} of an $m\times n$ matrix $A$ is
\[
\Nul A=\{\vect x\in\R^n:A\vect x=\vect0\}.
\]
It is a subspace of $\R^n$.
\pause

Indeed, $A\vect0=\vect0$, and if $A\vect u=A\vect v=\vect0$, then
\[
A(\vect u+\vect v)=A\vect u+A\vect v=\vect0.
\]
For any scalar $c$,
\[
A(c\vect u)=cA\vect u=\vect0.
\]
\end{frame}

\begin{frame}{Basis for a Subspace}
A \textbf{basis} for a subspace $H$ is a set of vectors that
\begin{enumerate}
  \item lies in $H$;
  \item spans $H$;
  \item is linearly independent.
\end{enumerate}
\pause

A basis contains enough vectors to generate the subspace, with no redundant vector.
\end{frame}

\begin{frame}{The Standard Basis}
The columns of $I_n$ form the standard basis for $\R^n$:
\[
\vect e_1=\begin{bmatrix}1\\0\\\vdots\\0\end{bmatrix},
\quad
\vect e_2=\begin{bmatrix}0\\1\\\vdots\\0\end{bmatrix},
\quad\ldots\quad,
\vect e_n=\begin{bmatrix}0\\0\\\vdots\\1\end{bmatrix}.
\]
Thus
\[
\{\vect e_1,\ldots,\vect e_n\}
\]
is a basis for $\R^n$.
\end{frame}

\begin{frame}{A Basis for a Null Space}
Let
\[
A=\begin{bmatrix}
-3&6&-1&1&-7\\
1&-2&2&3&-1\\
2&-4&5&8&-4
\end{bmatrix}.
\]
Row reduction of $A\vect x=\vect0$ gives
\[
\begin{bmatrix}
1&-2&0&-1&3\\
0&0&1&2&-2\\
0&0&0&0&0
\end{bmatrix}.
\]
The free variables are $x_2,x_4,x_5$.
\end{frame}

\begin{frame}{Null-Space Basis: Parametric Form}
The equations give
\[
x_1=2x_2+x_4-3x_5,
\qquad
x_3=-2x_4+2x_5.
\]
Hence
\[
\vect x=x_2\begin{bmatrix}2\\1\\0\\0\\0\end{bmatrix}
+x_4\begin{bmatrix}1\\0\\-2\\1\\0\end{bmatrix}
+x_5\begin{bmatrix}-3\\0\\2\\0\\1\end{bmatrix}.
\]
Therefore these three vectors form a basis for $\Nul A$.
\end{frame}

\begin{frame}{A Basis for a Column Space}
Row operations preserve the linear dependence relations among the columns, but they change the actual column vectors.
\pause

To find a basis for $\Col A$:
\begin{enumerate}
  \item row reduce $A$ to an echelon form;
  \item identify the pivot column indices;
  \item take those columns from the original matrix $A$.
\end{enumerate}
\pause
Do not use the pivot columns from the echelon form as vectors in $\Col A$.
\end{frame}

\begin{frame}{Pivot Columns and Column Space}
\theorembox{Pivot Column Theorem}{The pivot columns of a matrix $A$ form a basis for $\Col A$.}
\medskip
If the pivot columns of an echelon form occur in columns $1$, $2$, and $5$, then
\[
\{\vect a_1,\vect a_2,\vect a_5\}
\]
is a basis for $\Col A$, where the vectors come from the original matrix.
\end{frame}

% -----------------------------------------------------------------------------
\section[Dimension and Rank]{Dimension and Rank}
\sectioncard{2.9}{Dimension and Rank}

\begin{frame}{Why a Basis Matters}
Every vector in a subspace has a unique representation as a linear combination of basis vectors.
\pause

If $\beta=\{\vect b_1,\ldots,\vect b_p\}$ is a basis and
\[
\vect x=c_1\vect b_1+\cdots+c_p\vect b_p
=d_1\vect b_1+\cdots+d_p\vect b_p,
\]
then
\[
\vect0=(c_1-d_1)\vect b_1+\cdots+(c_p-d_p)\vect b_p.
\]
Linear independence forces $c_i=d_i$ for every $i$.
\end{frame}

\begin{frame}{Coordinate Vectors}
Let $\beta=\{\vect b_1,\ldots,\vect b_p\}$ be a basis for $H$.
If
\[
\vect x=c_1\vect b_1+\cdots+c_p\vect b_p,
\]
then the coordinate vector of $\vect x$ relative to $\beta$ is
\[
\boxed{[\vect x]_\beta=
\begin{bmatrix}c_1\\\vdots\\c_p\end{bmatrix}}.
\]
The basis determines the coordinate system.
\end{frame}

\begin{frame}{Coordinate Example}
Let
\[
\vect v_1=\begin{bmatrix}3\\6\\2\end{bmatrix},
\quad
\vect v_2=\begin{bmatrix}-1\\0\\1\end{bmatrix},
\quad
\vect x=\begin{bmatrix}3\\12\\7\end{bmatrix},
\quad
\beta=\{\vect v_1,\vect v_2\}.
\]
Solve $c_1\vect v_1+c_2\vect v_2=\vect x$:
\[
\left[\begin{array}{cc|c}3&-1&3\\6&0&12\\2&1&7\end{array}\right]
\sim
\left[\begin{array}{cc|c}1&0&2\\0&1&3\\0&0&0\end{array}\right].
\]
Thus
\[
\boxed{[\vect x]_\beta=\begin{bmatrix}2\\3\end{bmatrix}}.
\]
\end{frame}

\begin{frame}{Dimension}
The \textbf{dimension} of a nonzero subspace $H$, written $\dim H$, is the number of vectors in any basis for $H$.
\[
\dim\{\vect0\}=0.
\]
Examples:
\begin{itemize}
  \item $\dim\R^n=n$;
  \item a line through the origin has dimension $1$;
  \item a plane through the origin in $\R^3$ has dimension $2$.
\end{itemize}
\end{frame}

\begin{frame}{Rank and Nullity}
The \textbf{rank} of a matrix is the dimension of its column space:
\[
\boxed{\rank A=\dim(\Col A)}.
\]
The \textbf{nullity} is the dimension of its null space:
\[
\boxed{\operatorname{nullity}A=\dim(\Nul A)}.
\]
\pause
Rank equals the number of pivot columns. Nullity equals the number of free variables in $A\vect x=\vect0$.
\end{frame}

\begin{frame}{Rank Example}
Suppose row reduction gives
\[
A\sim
\begin{bmatrix}
2&5&-3&-4&8\\
0&-3&2&5&-7\\
0&0&0&4&-6\\
0&0&0&0&0
\end{bmatrix}.
\]
The pivot columns are $1$, $2$, and $4$. Therefore
\[
\boxed{\rank A=3}.
\]
If $A$ has five columns, then the two nonpivot variables are free, so $\dim(\Nul A)=2$.
\end{frame}

\begin{frame}{The Rank--Nullity Theorem}
\theorembox{Rank--Nullity Theorem}{If $A$ has $n$ columns, then}
\[
\boxed{\rank A+\dim(\Nul A)=n}.
\]
\pause
Every column corresponds either to
\begin{itemize}
  \item a pivot variable, counted by $\rank A$, or
  \item a free variable, counted by $\dim(\Nul A)$.
\end{itemize}
\end{frame}

\begin{frame}{A Basis Criterion in a Finite-Dimensional Space}
Let $H$ be a $p$-dimensional subspace of $\R^n$.
\begin{itemize}
  \item Any linearly independent set of exactly $p$ vectors in $H$ is automatically a basis for $H$.
  \item Any set of exactly $p$ vectors that spans $H$ is automatically a basis for $H$.
\end{itemize}
\pause
Once the number of vectors matches the dimension, either spanning or linear independence implies the other condition.
\end{frame}

\begin{frame}{The Invertible Matrix Theorem: New Forms}
For an $n\times n$ matrix $A$, the following are equivalent to invertibility:
\begin{itemize}
  \item the columns of $A$ form a basis for $\R^n$;
  \item $\Col A=\R^n$;
  \item $\dim(\Col A)=n$;
  \item $\rank A=n$;
  \item $\Nul A=\{\vect0\}$;
  \item $\dim(\Nul A)=0$.
\end{itemize}
\end{frame}

\begin{frame}{How the Equivalences Connect}
The logical chain is
\[
\Col A=\R^n
\Longrightarrow
\dim(\Col A)=n
\Longrightarrow
\rank A=n.
\]
Rank--nullity then gives
\[
\dim(\Nul A)=0
\Longrightarrow
\Nul A=\{\vect0\}.
\]
Thus $A\vect x=\vect0$ has only the trivial solution, one of the original forms of the Invertible Matrix Theorem.
\end{frame}

\begin{frame}{Lecture Summary}
\begin{itemize}
  \item Determinant properties make elimination and cofactor expansion efficient.
  \item A nonzero determinant is equivalent to invertibility for a square matrix.
  \item Rank comparison classifies all linear systems; Cramer's rule solves small invertible square systems.
  \item Column and null spaces are subspaces with bases obtained from pivots and free variables.
  \item Coordinates are unique relative to a basis, and rank--nullity divides the variables into pivot and free variables.
\end{itemize}
\end{frame}

\end{document}
